\chapter{Installation and Setup Guide}
\label{appendix:setup}

This appendix provides comprehensive instructions for installing, configuring, and running the Clinical RAG Chatbot system developed in this research project.

\section{System Requirements}

\subsection{Hardware Requirements}

\textbf{Minimum Requirements:}
\begin{itemize}
    \item \textbf{RAM}: 8GB (16GB recommended for optimal performance)
    \item \textbf{Storage}: 50GB free disk space for models and data
    \item \textbf{CPU}: Multi-core processor (4+ cores recommended)
    \item \textbf{GPU}: Optional but recommended for embedding model inference
\end{itemize}

\textbf{Recommended Specifications:}
\begin{itemize}
    \item \textbf{RAM}: 16GB or higher
    \item \textbf{Storage}: 100GB+ SSD storage
    \item \textbf{GPU}: NVIDIA GPU with 6GB+ VRAM (for faster model processing)
    \item \textbf{Network}: Stable internet connection for model downloads (5-10GB total)
\end{itemize}

\subsection{Software Requirements}

\textbf{Core Dependencies:}
\begin{itemize}
    \item \textbf{Python}: 3.10 or higher (3.11 recommended)
    \item \textbf{Node.js}: 16+ (v18+ recommended for frontend)
    \item \textbf{Git}: Latest version for repository management
    \item \textbf{Ollama}: For local LLM hosting and management
\end{itemize}

\textbf{Operating System Support:}
\begin{itemize}
    \item \textbf{Windows}: 10/11 (64-bit)
    \item \textbf{macOS}: 10.15+ (Catalina and newer)
    \item \textbf{Linux}: Ubuntu 20.04+, CentOS 8+, or equivalent distributions
\end{itemize}

\section{Installation Steps}

\subsection{Step 1: Install Prerequisites}

\textbf{Python Installation:}
\begin{itemize}
    \item Download Python 3.11 from \texttt{https://python.org}
    \item Ensure \texttt{pip} is included and \texttt{python} is added to PATH
    \item Verify installation: \texttt{python --version}
\end{itemize}

\textbf{Node.js Installation:}
\begin{itemize}
    \item Download Node.js 18+ from \texttt{https://nodejs.org}
    \item Verify installation: \texttt{node --version} and \texttt{npm --version}
\end{itemize}

\textbf{Ollama Installation:}
\begin{minted}[fontsize=\footnotesize, breaklines]{bash}
# Visit https://ollama.ai and follow platform-specific instructions
# For Windows: Download installer from website
# For macOS: brew install ollama
# For Linux: curl -fsSL https://ollama.ai/install.sh | sh

# Verify installation
ollama --version
\end{minted}

\subsection{Step 2: Clone Repository and Setup Environment}

\begin{minted}[fontsize=\footnotesize, breaklines]{bash}
# Clone the repository
git clone https://github.com/DevNick21/msc_project.git
cd msc_project

# Create Python virtual environment
python -m venv venv

# Activate virtual environment
# Windows:
venv\Scripts\activate

# macOS/Linux:
source venv/bin/activate
\end{minted}

\subsection{Step 3: Install Python Dependencies}

\begin{minted}[fontsize=\footnotesize, breaklines]{bash}
# Install project in development mode
pip install -e .

# Verify core dependencies
pip list | grep -E "langchain|torch|sentence-transformers"
\end{minted}

Key Python packages installed include:
\begin{itemize}
    \item \textbf{LangChain Framework}: v0.3.25 (Core RAG pipeline)
    \item \textbf{PyTorch}: v2.7.1 (Deep learning backend)
    \item \textbf{Sentence Transformers}: v4.1.0 (Embedding models)
    \item \textbf{FAISS}: v1.11.0 (Vector similarity search)
    \item \textbf{Flask}: v3.0.2 (API server)
    \item \textbf{Pandas}: v2.3.0 (Data processing)
    \item \textbf{Ollama}: v0.5.1 (LLM integration)
\end{itemize}

\subsection{Step 4: Install LLM Models via Ollama}

\begin{minted}[fontsize=\footnotesize, breaklines]{bash}
# Pull required LLM models (choose based on your needs)
ollama pull deepseek-r1:1.5b    # Reasoning-focused (default)
ollama pull qwen3:1.7b          # Multilingual capabilities
ollama pull llama3.2:1b         # Meta's efficient model
ollama pull gemma:2b            # Google's lightweight model
ollama pull phi3:3.8b           # Microsoft's instruction model
ollama pull tinyllama:1.1b      # Most compact model

# Verify models are installed
ollama list
\end{minted}

\textbf{Model Size Requirements:}
\begin{itemize}
    \item DeepSeek-R1 (1.5B): ~900MB
    \item Qwen3 (1.7B): ~1.0GB
    \item Llama3.2 (1B): ~670MB
    \item Gemma (2B): ~1.4GB
    \item Phi3 (3.8B): ~2.2GB
    \item TinyLlama (1.1B): ~600MB
\end{itemize}

\subsection{Step 5: Install Frontend Dependencies (Optional)}

\begin{minted}[fontsize=\footnotesize, breaklines]{bash}
# Navigate to frontend directory
cd frontend

# Install Node.js dependencies
npm install

# Verify installation
npm list --depth=0
\end{minted}

Key frontend packages include:
\begin{itemize}
    \item \textbf{React}: v18.2.0 (UI framework)
    \item \textbf{Material-UI}: v5.15.0 (Component library)
    \item \textbf{Emotion}: v11.11.0 (CSS-in-JS styling)
\end{itemize}

\section{Configuration}

\subsection{Model Configuration}

Edit \texttt{RAG\_chat\_pipeline/config/config.py} to customize model selection:

\begin{minted}[fontsize=\footnotesize, breaklines]{python}
# Default embedding model selection
model_in_use = "biomedbert"  # Options: biomedbert, mini-lm, 
                             # ms-marco, multi-qa, e5-base, 
                             # BioLORD, BioBERT, MedQuAD

# Default LLM model selection  
LLM_MODEL = llms["deepseek"]  # Options: deepseek, qwen, llama,
                              # gemma, phi3, tinyllama
\end{minted}

\textbf{Available Embedding Models:}
\begin{itemize}
    \item \textbf{biomedbert}: Medical domain-optimized (recommended for clinical data)
    \item \textbf{multi-qa}: Question-answering specialized
    \item \textbf{mini-lm}: Lightweight general-purpose
    \item \textbf{ms-marco}: Search-optimized
    \item \textbf{e5-base}: Efficient general embeddings
    \item \textbf{BioLORD}: Biomedical language representation
    \item \textbf{BioBERT}: Biomedical BERT variant
    \item \textbf{MedQuAD}: Medical QA specialized
\end{itemize}

\subsection{System Performance Configuration}

Key performance parameters in \texttt{config.py}:

\begin{minted}[fontsize=\footnotesize, breaklines]{python}
# Retrieval settings
DEFAULT_K = 5                    # Documents retrieved per query
RETRIEVAL_MAX_K = 5              # Max docs for focused searches
GLOBAL_SEARCH_MAX_K = 20         # Max docs for global searches
CANDIDATE_DOC_LIMIT = 20         # Max candidates after filtering

# Chat history management
MAX_CHAT_HISTORY = 60            # Max conversation messages

# Feature flags for optimization
ENABLE_REPHRASING = False        # Query rephrasing (slower)
ENABLE_ENTITY_EXTRACTION = False # Auto parameter extraction
LOG_LEVEL = "quiet"              # Logging verbosity
\end{minted}

\subsection{Data Configuration}

The system automatically selects between real MIMIC-IV data and synthetic data:

\textbf{Synthetic Data (Default):}
\begin{itemize}
    \item Automatically generated if real MIMIC-IV data unavailable
    \item Contains 100 fictional patients with 150 admissions
    \item Includes all clinical data types (diagnoses, labs, procedures, etc.)
    \item Safe for public distribution and demonstration
\end{itemize}

\textbf{Real MIMIC-IV Data (Optional):}
\begin{itemize}
    \item Requires credentialed access to MIMIC-IV dataset
    \item Place sample data in \texttt{mimic\_sample\_1000/} directory
    \item Run data processing notebooks in \texttt{data\_handling/}
\end{itemize}

\section{Running the System}

\subsection{Automated Startup (Recommended)}

\textbf{Windows Users:}
\begin{minted}[fontsize=\footnotesize, breaklines]{bash}
# Double-click or execute:
start_app.bat
\end{minted}

\textbf{macOS/Linux Users:}
\begin{minted}[fontsize=\footnotesize, breaklines]{bash}
# Make executable and run:
chmod +x start_app.sh
./start_app.sh
\end{minted}

This automated script will:
\begin{enumerate}
    \item Create virtual environment if needed
    \item Activate virtual environment
    \item Install/update dependencies
    \item Start Flask API server (port 5000)
    \item Start React frontend (port 3000)
    \item Open web browser automatically
\end{enumerate}

\subsection{Manual Startup}

\textbf{Backend API Server:}
\begin{minted}[fontsize=\footnotesize, breaklines]{bash}
# Ensure virtual environment is activated
# Windows: venv\Scripts\activate
# macOS/Linux: source venv/bin/activate

# Start Flask API server
python api/app.py

# Server will start on http://localhost:5000
\end{minted}

\textbf{Frontend Web Interface:}
\begin{minted}[fontsize=\footnotesize, breaklines]{bash}
# In a new terminal, navigate to frontend directory
cd frontend

# Start React development server
npm start

# Interface will open at http://localhost:3000
\end{minted}

\textbf{Command Line Interface:}
\begin{minted}[fontsize=\footnotesize, breaklines]{bash}
# Interactive CLI chat interface
python cli_chat.py chat

# Direct query mode
python -m RAG_chat_pipeline.core.main
\end{minted}

\subsection{System Evaluation and Testing}

\textbf{Quick System Test:}
\begin{minted}[fontsize=\footnotesize, breaklines]{bash}
# Test with specific model combination
python -m RAG_chat_pipeline.benchmarks.model_evaluation_runner \
    single mini-lm deepseek --type quick

# Quick evaluation with 3 questions
python -m RAG_chat_pipeline.benchmarks.rag_evaluator quick
\end{minted}

\textbf{Short Evaluation:}
\begin{minted}[fontsize=\footnotesize, breaklines]{bash}
# Evaluate specific model combination (5 questions)
python -m RAG_chat_pipeline.benchmarks.model_evaluation_runner \
    single biomedbert deepseek --type short

# Short evaluation across all models
python -m RAG_chat_pipeline.benchmarks.model_evaluation_runner \
    all --type short
\end{minted}

\textbf{Full Model Comparison:}
\begin{minted}[fontsize=\footnotesize, breaklines]{bash}
# Complete evaluation of all 54 model combinations
python -m RAG_chat_pipeline.benchmarks.model_evaluation_runner \
    all --type full

# Generate comprehensive report
python -m RAG_chat_pipeline.benchmarks.model_evaluation_runner report
\end{minted}

\section{Accessing the System}

\subsection{Web Interface Usage}

Navigate to \texttt{http://localhost:3000} after startup:

\begin{enumerate}
    \item \textbf{Chat Interface}: Type natural language questions about clinical data
    \item \textbf{Model Selection}: Choose embedding and LLM model combinations
    \item \textbf{Session Management}: View conversation history and clear sessions
    \item \textbf{Sample Questions}: Use provided examples to explore functionality
\end{enumerate}

\textbf{Example Queries:}
\begin{itemize}
    \item "What diagnoses does admission 25282710 have?"
    \item "Show abnormal lab values for patient 12345"
    \item "What medications was the patient prescribed?"
    \item "Were any cultures positive for infections?"
\end{itemize}

\subsection{API Endpoints}

The Flask API provides programmatic access:

\textbf{Chat Endpoint:}
\begin{minted}[fontsize=\footnotesize, breaklines]{bash}
POST http://localhost:5000/api/chat
Content-Type: application/json

{
  "query": "What lab values were abnormal?",
  "hadm_id": 25282710,
  "history": []
}
\end{minted}

\textbf{Models Endpoint:}
\begin{minted}[fontsize=\footnotesize, breaklines]{bash}
GET http://localhost:5000/api/models

# Returns available embedding models and vector stores
\end{minted}

\section{Troubleshooting}

\subsection{Common Issues and Solutions}

\textbf{Python Environment Issues:}
\begin{itemize}
    \item \textbf{Issue}: Import errors or module not found
    \item \textbf{Solution}: Ensure virtual environment is activated and run \texttt{pip install -e .}
\end{itemize}

\textbf{Ollama Model Issues:}
\begin{itemize}
    \item \textbf{Issue}: LLM model not responding or connection errors
    \item \textbf{Solution}: Verify Ollama is running (\texttt{ollama serve}) and models are pulled
\end{itemize}

\textbf{Port Conflicts:}
\begin{itemize}
    \item \textbf{Issue}: Port 5000 or 3000 already in use
    \item \textbf{Solution}: Modify ports in \texttt{api/app.py} and \texttt{frontend/package.json}
\end{itemize}

\textbf{Memory Issues:}
\begin{itemize}
    \item \textbf{Issue}: Out of memory during model loading
    \item \textbf{Solution}: Use smaller models (tinyllama, mini-lm) or increase system RAM
\end{itemize}

\textbf{Frontend Build Issues:}
\begin{itemize}
    \item \textbf{Issue}: npm install fails or React won't start
    \item \textbf{Solution}: Clear npm cache (\texttt{npm cache clean --force}) and reinstall
\end{itemize}

\subsection{Performance Optimization}

\textbf{For Low-Resource Systems:}
\begin{itemize}
    \item Use \texttt{tinyllama} or \texttt{mini-lm} models
    \item Reduce \texttt{DEFAULT\_K} to 3 in configuration
    \item Set \texttt{LOG\_LEVEL} to "quiet"
    \item Disable entity extraction and rephrasing features
\end{itemize}

\textbf{For High-Performance Systems:}
\begin{itemize}
    \item Use \texttt{phi3} or \texttt{biomedbert} models
    \item Increase \texttt{DEFAULT\_K} to 10 for better recall
    \item Enable GPU acceleration for PyTorch
    \item Use multiple embedding models simultaneously
\end{itemize}

\section{Data Management}

\subsection{Synthetic Data Generation}

\begin{minted}[fontsize=\footnotesize, breaklines]{bash}
# Generate fresh synthetic data
python -m synthetic_data.synthetic_data_generator

# Process synthetic data into vector stores
# (Run notebooks in data_handling/)
\end{minted}

\subsection{MIMIC-IV Data Processing}

If using real MIMIC-IV data:
\begin{enumerate}
    \item Place CSV files in \texttt{mimic\_sample\_1000/}
    \item Run \texttt{data\_handling/creating\_docs.ipynb}
    \item Execute all cells to create chunked documents
    \item Vector stores will be automatically generated
\end{enumerate}

\section{System Validation}

\subsection{Installation Verification Checklist}

\begin{itemize}
    \item[$\square$] Python 3.10+ installed and accessible
    \item[$\square$] Virtual environment created and activated
    \item[$\square$] All pip dependencies installed without errors
    \item[$\square$] Ollama installed and running
    \item[$\square$] At least one LLM model pulled via Ollama
    \item[$\square$] Node.js and npm installed (if using frontend)
    \item[$\square$] Frontend dependencies installed (if using web interface)
    \item[$\square$] Flask API server starts without errors
    \item[$\square$] React frontend compiles and serves (if using web interface)
    \item[$\square$] Quick evaluation test completes successfully
\end{itemize}

\subsection{Final Verification Commands}

\begin{minted}[fontsize=\footnotesize, breaklines]{bash}
# Verify Python environment
python --version
pip list | grep langchain

# Verify Ollama
ollama list

# Verify system functionality
python -m RAG_chat_pipeline.benchmarks.rag_evaluator quick

# Test API endpoint (if server running)
curl http://localhost:5000/api/models
\end{minted}

Upon successful completion of this setup guide, the Clinical RAG system will be fully operational and ready for clinical data querying, evaluation, and research applications as described in this dissertation.
